\documentclass{ximera}

\title{Project 3: Jose Miguel}
\author{Jose Miguel}

\begin{document}
\begin{abstract}
    Project 3 on kaleidoscope completed by Jose Miguel .
\end{abstract}

\maketitle

\begin{problem}
   Find the linear transformation that reflects an image by a mirror on the x-axis.
   
   \vspace{0.3cm}
   
  A reflection across the x-axis maps a point $(x, y)$ to $(x, -y)$.


     \vspace{0.3cm}

This transformation can be written in matrix form as:

    \[
        \begin{bmatrix}
           1 & 0 \\
           0 & -1
        \end{bmatrix}
        \begin{bmatrix}
           x \\
           y
        \end{bmatrix}
           =
        \begin{bmatrix}
           x \\
           - y
        \end{bmatrix}
    \]

     \vspace{0.3cm}

Thus, the matrix representing reflection across the x-axis is:

\[
        \begin{bmatrix}
           1 & 0 \\
           0 & -1
        \end{bmatrix}
\]
        \end{problem}

    \vspace{0.5cm}

        \begin{problem}
Find the linear transformation that rotates an image counterclockwise through an angle of $60^\circ$.

    \vspace{0.3cm}

    A rotation by an angle $\theta$ transforms vectors in the plane using a rotation matrix.

     The standard rotation matrix is given by:

\[
\begin{bmatrix}
\cos\theta & -\sin\theta \\
\sin\theta & \cos\theta
\end{bmatrix}
\]

\vspace{0.3cm}

For a rotation of 60 degrees, we substitute $\theta = 60^\circ$:

\[
\begin{bmatrix}
\cos 60^\circ & -\sin 60^\circ \\
\sin 60^\circ & \cos 60^\circ
\end{bmatrix}
\]

\vspace{0.3cm}

Since $\cos 60^\circ = \frac{1}{2}$ and $\sin 60^\circ = \frac{\sqrt{3}}{2}$, we obtain:

\[
\begin{bmatrix}
\frac{1}{2} & -\frac{\sqrt{3}}{2} \\
\frac{\sqrt{3}}{2} & \frac{1}{2}
\end{bmatrix}
\]

\end{problem}

       \vspace{0.5cm}

    \begin{problem}
    Find the linear transformation that reflects an image by a mirror on the diagonal line that makes an angle of $60^\circ$ with the x-axis

     \vspace{0.3cm}

    The reflection matrix across a line making an angle $\theta$ with the x-axis is given by:

\[
\begin{bmatrix}
\cos(2\theta) & \sin(2\theta) \\
\sin(2\theta) & -\cos(2\theta)
\end{bmatrix}
\]

   \vspace{0.3cm}
   
For a line making an angle of $60^\circ$, we substitute $\theta = 60^\circ$:

\[
\begin{bmatrix}
\cos(120^\circ) & \sin(120^\circ) \\
\sin(120^\circ) & -\cos(120^\circ)
\end{bmatrix}
\]

   \vspace{0.3cm}

Since $\cos 120^\circ = -\frac{1}{2}$ and $\sin120^\circ = \frac{\sqrt{3}}{2}$, we obtain:

\[
\begin{bmatrix}
-\frac{1}{2} & \frac{\sqrt{3}}{2} \\
\frac{\sqrt{3}}{2} & \frac{1}{2}
\end{bmatrix}
\]

  \vspace{0.3cm}
  
This matrix was developed using the general formula for reflection across a line at an angle $\theta$, which involves doubling the angle to account for the geometry of reflection.
    \end{problem}

    \vspace{0.6cm}

\begin{center}
    \textbf{Project Extension: The Algebraic Structure of the Kaleidoscope\\}
\end{center}

\begin{problem} 
\textbf{Problem 4a.}\\
  Here, we know R is the rotation matrix by $60^\circ$ from "Problem 2" and F is the reflection matrix across the x-axis.\\
  Hence:

 $$ R =
\begin{bmatrix}
\frac{1}{2} & -\frac{\sqrt{3}}{2} \\
\frac{\sqrt{3}}{2} & \frac{1}{2}
\end{bmatrix} \quad 
F =
\begin{bmatrix}
1 & 0 \\
0 & -1
\end{bmatrix}$$
  
  (a) Computing $RF$ and $FR$:\\
  
  \vspace{0.04cm}
  
    First, we compute $RF$:

  \vspace{0.3cm}

$$RF =
\begin{bmatrix}
\frac{1}{2} & -\frac{\sqrt{3}}{2} \\
\frac{\sqrt{3}}{2} & \frac{1}{2}
\end{bmatrix}
\begin{bmatrix}
1 & 0 \\
0 & -1
\end{bmatrix}
=
\begin{bmatrix}
\frac{1}{2} & \frac{\sqrt{3}}{2} \\
\frac{\sqrt{3}}{2} & -\frac{1}{2}
\end{bmatrix}$$

Next, we compute $FR$:

$$FR =
\begin{bmatrix}
1 & 0 \\
0 & -1
\end{bmatrix}
\begin{bmatrix}
\frac{1}{2} & -\frac{\sqrt{3}}{2} \\
\frac{\sqrt{3}}{2} & \frac{1}{2}
\end{bmatrix}
=
\begin{bmatrix}
\frac{1}{2} & -\frac{\sqrt{3}}{2} \\
-\frac{\sqrt{3}}{2} & -\frac{1}{2}
\end{bmatrix}$$

(b) Since $RF \neq FR$, the matrices do not commute.

\vspace{0.3cm}

(c) My calculation shows that the order of transformations matters. Applying rotation first and then reflection produces a different result than applying reflection first and then rotation.

\vspace{0.3cm}

(d) Visually, applying reflection first and then rotation produces a different image than applying rotation first and then reflection. The shape will end up in a different position and orientation, showing that the order of transformations changes the final result.

  \vspace{0.7cm}

\textbf{Problem 4b.}\\

  \vspace{0.01cm}
  
(a) Here we have to compute $R^2$, $R^3$, and $R^6$\\

    \vspace{0.03cm}
    
   $$ R^2(120^\circ) =
\begin{bmatrix}
\cos 120^\circ & -\sin 120^\circ \\
\sin 120^\circ & \cos 120^\circ
\end{bmatrix}
=
\begin{bmatrix}
-\frac{1}{2} & -\frac{\sqrt{3}}{2} \\
\frac{\sqrt{3}}{2} & -\frac{1}{2}
\end{bmatrix} $$

  \vspace{0.2cm}

$$ R^3(180^\circ) =
\begin{bmatrix}
\cos 180^\circ & -\sin 180^\circ \\
\sin 180^\circ & \cos 180^\circ
\end{bmatrix}
=
\begin{bmatrix}
-1 & 0 \\
0 & -1
\end{bmatrix}$$

  \vspace{0.2cm}

 $$ R^6(360^\circ) =
\begin{bmatrix}
\cos 360^\circ & -\sin 360^\circ \\
\sin 360^\circ & \cos 360^\circ
\end{bmatrix}
=
\begin{bmatrix}
1 & 0 \\
0 & 1
\end{bmatrix}$$

  \vspace{0.3cm}

(b) From the pattern above, we observe that each power of $R$ increases the angle by $60^\circ$. Making a conjecture about $R^k$ for any positive integer $k$

  \vspace{0.3cm}

$$  R^k =
\begin{bmatrix}
\cos(60k^\circ) & -\sin(60k^\circ) \\
\sin(60k^\circ) & \cos(60k^\circ)
\end{bmatrix}$$

  \vspace{0.3cm}

(c) Since $R$ represents a rotation of $60^\circ$, applying it six times results in a total rotation of $360^\circ$, which returns all points to their original position. Therefore, $R^6 = I$.

   \vspace{0.7cm}

\textbf{Problem 4c.}\\
  
We compute the determinants of the given matrices.

First, for the rotation matrix $R$:

\[
\det(R) =
\left(\frac{1}{2} \cdot \frac{1}{2}\right)
-
\left(-\frac{\sqrt{3}}{2} \cdot \frac{\sqrt{3}}{2}\right)
= \frac{1}{4} + \frac{3}{4} = 1
\]

Next, for the reflection matrix $F$:

\[
\det(F) = (1)(-1) - (0)(0) = -1
\]

For the composition $RF$, we use the property of determinants:

\[
\det(RF) = \det(R)\det(F) = (1)(-1) = -1
\]

  \vspace{0.3cm}

The determinant shows how a transformation changes a shape. A value of 1 means the shape keeps the same size and orientation, while a value of -1 means the size stays the same but the shape is flipped.

Thus, all kaleidoscope transformations preserve distances (and therefore area), but reflections reverse orientation while rotations do not.

\end{problem}
\end{document}
